% Werke X/1 (1917), native210–213, printed203–206.
\begingroup
\setlength{\headheight}{25pt}
\makeatletter\tagsleft@true\let\veqno\@@leqno\makeatother
\renewcommand{\thefootnote}{[*)}
\SourceLeaf{210}\ArithmeticHead{203}{R}{}
\fancyhead[C]{\scriptsize THEORIA SINUS LEMNISCATICI UNIVERSALISSIME ACCEPTI.}
\[\frac{(\Pi z^4)^6}{(\Pi zz)^4}=\frac1{M\cos\zeta};\quad
\frac{(\Pi y^4)^6}{(\Pi yy)^4}=\frac1{M\cos\eta}\quad\text{[*)]}\]
\footnotetext{In der Handschrift steht links vom Gleichheitszeichen beidemal \(M\) statt \(\Pi\).]}
\[\boxed{\begin{gathered}
\frac{(\Pi y^2)^6(\Pi z)^4}{(\Pi z^2)^6(\Pi y)^4}=-\frac{\log z}\pi=-\frac\pi{\log y}\\[.7em]
\log y\log z=\pi\pi
\end{gathered}}\]
\NumberedItem{11}
\noindent [S.\ 42]
{\small
\[\begin{aligned}
(P\varphi)^2&=0{,}3522376226\;6118372314&&=A\\
&-0{,}3535519576\;3585935635\cos2\varphi&&=2B[e^{-\frac12\pi}]\\
&+0{,}0013155679\;2352259042\cos4\varphi&&=2Ae^{-2\pi}\\
&-0{,}0000012329\;5741446398\cos6\varphi&&=2B[e^{-\frac12\pi}]e^{-4\pi}\\
&+0{,}.\,.\,.\,.\,.\,.\,.\,.\,.\,.\;0856752170\cos8\varphi&&=2Ae^{-8\pi}\\
&-0{,}.\,.\,.\,.\,.\,.\,.\,.\,.\,.\; .\,.\,.\,.\,.\,.\;1494\cos10\varphi&&=2B[e^{-\frac12\pi}]e^{-12\pi}
\end{aligned}\]}
\[\begin{aligned}
(P\varphi)^2={}&A(1+2e^{-2\pi}\cos4\varphi+2e^{-8\pi}\cos8\varphi+2e^{-18\pi}\cos12\varphi\text{ etc.})\\
&-2B(e^{-\frac12\pi}\cos2\varphi+e^{-\frac92\pi}\cos6\varphi+e^{-\frac{25}2\pi}\cos10\varphi\text{ etc.})
\end{aligned}\]
\[B=(1+\surd2)A,\quad
A=\frac{\surd\frac\pi\varpi}{\cos22\frac12{}^0\surd8[{}^4\!\surd2]}.\]
\noindent [S.\ 43]
\[0=\frac{2^{\frac14}\cos0{,}25}{\surd\frac\pi\varpi}A-\frac{2^{\frac14}\sin0{,}25}{\surd\frac\pi\varpi}B\]
\[\surd\frac12=\frac{2^{\frac14}\cos0{,}25}{\surd\frac\pi\varpi}A+\frac{2^{\frac14}\sin0{,}25}{\surd\frac\pi\varpi}B\]
\[A=2^{-\frac74}\frac{\surd\frac\pi\varpi}{\cos0{,}25},\quad
B=2^{-\frac74}\frac{\surd\frac\pi\varpi}{\sin0{,}25}.\]
\par\hfill26\textsuperscript{*}
\clearpage
\SourceLeaf{211}\ArithmeticHead{204}{L}{ANALYSIS.\quad NACHLASS.}
\NumberedItem{12}
\par\indent Si
\[\begin{aligned}
&(1+ax)(1+ax^3)(1+ax^5)\ldots\left(1+\frac xa\right)\left(1+\frac{x^3}a\right)\left(1+\frac{x^5}a\right)\ldots=K
\end{aligned}\]
\noindent supponitur producere
\[+\frac R{aa}+\frac Qa+P+Qa+Raa+\ldots,\]
\noindent producet
{\small
\[\begin{gathered}
(1+ax^3)(1+ax^5)\ldots\left(1+\frac1{ax}\right)\left(1+\frac xa\right)\ldots
=K\times\frac{1+\frac1{ax}}{1+ax}\;\text{[*)]}\\
\ldots+\frac R{x^4}\frac1{aa}+\frac Q{xx}\frac1a+P+Qxxa+Rx^4aa+\ldots
\end{gathered}\]}
\footnotetext{In der Handschrift stehen links vom Gleichheitszeichen statt der positiven Vorzeichen negative.]}
\noindent Quare]
\[\begin{aligned}
[K]&=\ldots\frac R{x^3}.\frac1a+\frac Qx+Pxa+Qx^3aa+\text{etc.}\\
&=\ldots Q\left[\frac1a+\right]P[+]Q[a+]R[aa+\ldots]\\
&=P\left(1+x\left(a+\frac1a\right)+x^4\left(aa+\frac1{aa}\right)+x^9\left(a^3+\frac1{a^3}\right)\ldots\right).
\end{aligned}\]
\ClosingRule
\par\indent Si ex quadrato [ipsius \(K\)] prodit
\[+\frac Qa+P+Qa\text{ etc.,}\]
\noindent erit
{\small
\[\begin{gathered}
[\ldots+]\frac P{xxaa}[+]\frac Q{xxa}[+]\frac R{xx}[+]\frac S{xx}a[+]\frac T{xx}[aa+\ldots]\\
{}[=\ldots+]\frac R{x^4aa}[+]\frac Q{xxa}[+]P[+]Qxxa[+]Rx^4[aa+\ldots]
\end{gathered}\]}
\noindent [Quare]
{\small
\[\begin{aligned}
[K^2={}]P[+]&Q\left[\left(a+\frac1a\right)+\right]Pxx\left[\left(aa+\frac1{aa}\right)+\right]
Qx^4\left[\left(a^3+\frac1{a^3}\right)+\right]\\
&Px^8\left[\left(a^4+\frac1{a^4}\right)+\right]Qx^{12}\left[\left(a^5+\frac1{a^5}\right)+\right]
Px^{18}\left[\left(a^6+\frac1{a^6}\right)+\text{etc.}\right]
\end{aligned}\]}
\clearpage
\SourceLeaf{212}\ArithmeticHead{205}{R}{}
\fancyhead[C]{\scriptsize THEORIA SINUS LEMNISCATICI UNIVERSALISSIME ACCEPTI.}
\NumberedItem{13}
\noindent [S.\ 45]
\[\begin{aligned}
T&=0&\varphi&=k\varpi+\tfrac12l\varpi'\surd-1\\
W&=0&\varphi&=(k+\tfrac12)\varpi+(l+\tfrac12)\varpi\surd-1
\end{aligned}\]
{\small
\[\begin{aligned}
Tn\varpi={}&\frac\varpi\pi\sin n\pi
\left(1+\frac{4\sin n\pi^2}{\left(e^{\frac{\varpi'}\varpi\pi}-e^{-\frac{\varpi'}\varpi\pi}\right)^2}\right)\\
&\quad\left(1+\frac{4\sin n\pi^2}{\left(e^{2\frac{\varpi'}\varpi\pi}-e^{-2\frac{\varpi'}\varpi\pi}\right)^2}\right)\ldots
\end{aligned}\]}
\[Qn\varpi=\left(1-\frac{4\sin n\pi}{\left(e^{\frac12\frac{\varpi'}\varpi\pi}+e^{\frac12\frac{\varpi'}\varpi\pi}\right)^2}\right)\ldots\]
\ClosingRule
\[\begin{aligned}
&\sin\varphi\left(1-\frac{2\cos\varphi}{x+x^{-1}}\right)\left(1+\frac{2\cos\varphi}{x+x^{-1}}\right)
\left(1-\frac{2\cos\varphi}{xx+x^{-2}}\right)(1+\ldots)\text{ etc.}\\
&\qquad=\frac{\dfrac{\sin\varphi}{x^{\frac14}}-\dfrac{\sin3\varphi}{x^{\frac94}}+\dfrac{\sin5\varphi}{x^{\frac{25}4}}\ldots}
{\dfrac1{x^{\frac14}}+\dfrac1{x^{\frac94}}+\dfrac1{x^{\frac{25}4}}\ldots}
\end{aligned}\]
\noindent [S.\ 46]
\[\begin{aligned}
&\left(1-\frac{2\cos2\varphi}{xx+\frac1{xx}}\right)
\left(1-\frac{2\cos2\varphi}{x^6+\frac1{x^6}}\right)
\left(1-\frac{2\cos2\varphi}{x^{10}+\frac1{x^{10}}}\right)\text{ etc.}\\
&\qquad=\frac{1-\dfrac{2\cos2\varphi}{xx}+\dfrac{2\cos4\varphi}{x^8}-\dfrac{2\cos6\varphi}{x^{18}}+\text{etc.}}
{1+\dfrac2{x^8}+\dfrac2{x^{32}}+\text{etc.}}
\end{aligned}\]
\[\frac{x^4+1}{x^4+2xx+1}.\frac{x^{12}+1}{x^{12}+2x^6+1}\ldots
=\frac{1-\dfrac2{x^8}+\dfrac2{x^{32}}-\text{etc.}}{1+\dfrac2{x^2}+\dfrac2{x^8}+\text{etc.}}\]
\ClosingRule
{\small
\[\begin{aligned}
&1+2y+2y^4+2y^9+\ldots\\
&\qquad=\frac{(1+y)^2(1+y^3)^2(1+y^4)^2(1+y^5)^2(1+y^7)^2\ldots}
{(1+yy)(1+y^6)(1+y^8)(1+y^{10})(1+y^{14})\ldots}.
\end{aligned}\]}
\clearpage
\SourceLeaf{213}\ArithmeticHead{206}{L}{ANALYSIS.\quad NACHLASS.}
\NumberedItem{14}
\noindent [S.\ 51]
\[x=e^{\frac12\frac{\varpi\pi}{\varpi'}};\quad a=\cos\varphi+\varepsilon\sin\varphi\]
{\small
\[T\varphi=\frac\varpi\pi\sin\varphi\,
\frac{\begin{gathered}(1-x^4a^2)(1-x^8a^2)(1-x^{12}a^2)\ldots\\
\left(1-\frac{x^4}{a^2}\right)\left(1-\frac{x^8}{a^2}\right)\left(1-\frac{x^{12}}{a^2}\right)\text{ etc.}\end{gathered}}
{(1-x^4)^2(1-x^8)^2(1-x^{12})^2\text{ etc.}}\]
\[W\varphi=\frac{\begin{gathered}(1+xxaa)(1+x^6aa)(1+x^{10}aa)\ldots\\
\left(1+\frac{x^2}{aa}\right)\left(1+\frac{x^6}{a^3}\right)\left(1+\frac{x^8}{a^3}\right)\text{ etc.}\end{gathered}}
{(1+x^2)^2(1+x^6)^2(1+x^{10})^2\text{ etc.}}.\]
}
\ClosingRule
\begin{center}
Functiones Lemniscaticas considerare\\
cöperamus 1797. Januar. 8.
\end{center}
\par\bigskip\centerline{\rule{.35\linewidth}{.4pt}}
\clearpage\endgroup
